\clearpage
\zlabel{triptych:concise:chronology:start}
\section{Scriptural Date and Location}
% Generated. Do not edit.
% Deterministic projection of research/chronology.toml.
% schema: 3
% document: liturgy/roman-rite/1962/propers/temporal/59-nineteenth-after-pentecost
% generated-by: tools/tpt proper-chronology annotations
\newcommand{\chronologyannotationclaim}[6]{#6}
\newcommand{\chronologyannotationcomparisonclaim}[8]{#8}
\newcommand{\chronologyannotationreach}[3]{}
\newcommand{\chronologyannotationgroup}[3]{#3}
\newcommand{\chronologyannotationcomparisongroup}[4]{#4}
\newcommand{\chronologyannotation}[1]{%
  \ifcsname triptychchronologyannotation@#1\endcsname
    \csname triptychchronologyannotation@#1\endcsname
  \else
    \PackageError{triptych}{No chronology annotation for #1}{}%
  \fi}
% comparison-dependency: src/gpt/liturgy/roman-rite/1962/propers/temporal/59-nineteenth-after-pentecost/research/chronology-profile-comparisons.toml sha256=7306dbbfa42399c1dab06cac21ca7fbdbf67fd0cb1fef34297adddca32b80cbc
% introit: Introit
\expandafter\def\csname triptychchronologyannotation@introit\endcsname{%
  \chronologyannotationgroup{composition}{preferred}{\textbf{Composition}: \chronologyannotationreach{Ps.77.1}{Ps}{true}\chronologyannotationclaim{critical.psalms.latest-composition-boundary}{composition}{catholic-critical-v1}{preferred}{before the Maccabean period, around 165 B.C.; no individual psalm can be dated securely}{Before c. 165 B.C.}}%
}
% epistle: Epistle
\expandafter\def\csname triptychchronologyannotation@epistle\endcsname{%
  \chronologyannotationgroup{composition}{disputed}{\textbf{Composition} -- disputed: \chronologyannotationreach{Eph.4.23}{Eph}{true}\chronologyannotationreach{Eph.4.24}{Eph}{true}\chronologyannotationreach{Eph.4.25}{Eph}{true}\chronologyannotationreach{Eph.4.26}{Eph}{true}\chronologyannotationreach{Eph.4.27}{Eph}{true}\chronologyannotationreach{Eph.4.28}{Eph}{true}\chronologyannotationclaim{composition.epistle-to-the-ephesians}{composition}{catholic-traditional-v1}{disputed}{a period between 58 and 63}{A.D. 58--63}; \chronologyannotationreach{Eph.4.23}{Eph}{true}\chronologyannotationreach{Eph.4.24}{Eph}{true}\chronologyannotationreach{Eph.4.25}{Eph}{true}\chronologyannotationreach{Eph.4.26}{Eph}{true}\chronologyannotationreach{Eph.4.27}{Eph}{true}\chronologyannotationreach{Eph.4.28}{Eph}{true}\chronologyannotationclaim{composition.epistle-to-the-ephesians}{composition}{catholic-traditional-v1}{disputed}{(Philemon; Colossians; Ephesians; Philippians), 61}{A.D. 61}.}%
  \space \chronologyannotationcomparisongroup{catholic-critical-v1}{composition}{composition-only}{\textbf{Comparison (catholic-critical-v1) -- Composition} -- composition-only: \chronologyannotationreach{Eph.4.23}{Eph}{true}\chronologyannotationreach{Eph.4.24}{Eph}{true}\chronologyannotationreach{Eph.4.25}{Eph}{true}\chronologyannotationreach{Eph.4.26}{Eph}{true}\chronologyannotationreach{Eph.4.27}{Eph}{true}\chronologyannotationreach{Eph.4.28}{Eph}{true}\chronologyannotationcomparisonclaim{critical.ephesians-later-disciple}{composition}{catholic-critical-v1}{catholic-critical-v1}{disputed}{passage.united-states-conference-of-catholic-bishops.new-american-bible-revised-edition.english-usccb-web-2026-09-28.ephesians-introduction}{around A.D. 80--100}{Ephesians under the later-disciple hypothesis (approximate), A.D. 80--100}.}%
}
% gradual: Gradual
\expandafter\def\csname triptychchronologyannotation@gradual\endcsname{%
  \chronologyannotationgroup{composition}{preferred}{\textbf{Composition}: \chronologyannotationreach{Ps.140.2}{Ps}{true}\chronologyannotationclaim{critical.psalms.latest-composition-boundary}{composition}{catholic-critical-v1}{preferred}{before the Maccabean period, around 165 B.C.; no individual psalm can be dated securely}{Before c. 165 B.C.}}%
}
% alleluia: Alleluia
\expandafter\def\csname triptychchronologyannotation@alleluia\endcsname{%
  \chronologyannotationgroup{composition}{preferred}{\textbf{Composition}: \chronologyannotationreach{Ps.104.1}{Ps}{true}\chronologyannotationclaim{critical.psalms.latest-composition-boundary}{composition}{catholic-critical-v1}{preferred}{before the Maccabean period, around 165 B.C.; no individual psalm can be dated securely}{Before c. 165 B.C.}}%
}
% gospel: Gospel
\expandafter\def\csname triptychchronologyannotation@gospel\endcsname{%
  \chronologyannotationgroup{narrated-event}{research-pending}{\textbf{Event} -- Narrated event date unresolved.}%
  \space \chronologyannotationgroup{composition}{disputed}{\textbf{Composition} -- disputed: \chronologyannotationreach{Matt.22.1}{Matt}{true}\chronologyannotationreach{Matt.22.10}{Matt}{true}\chronologyannotationreach{Matt.22.11}{Matt}{true}\chronologyannotationreach{Matt.22.12}{Matt}{true}\chronologyannotationreach{Matt.22.13}{Matt}{true}\chronologyannotationreach{Matt.22.14}{Matt}{true}\chronologyannotationreach{Matt.22.2}{Matt}{true}\chronologyannotationreach{Matt.22.3}{Matt}{true}\chronologyannotationreach{Matt.22.4}{Matt}{true}\chronologyannotationreach{Matt.22.5}{Matt}{true}\chronologyannotationreach{Matt.22.6}{Matt}{true}\chronologyannotationreach{Matt.22.7}{Matt}{true}\chronologyannotationreach{Matt.22.8}{Matt}{true}\chronologyannotationreach{Matt.22.9}{Matt}{true}\chronologyannotationclaim{composition.gospel-of-matthew}{composition}{catholic-traditional-v1}{disputed}{about A.D. 38-45}{c. A.D. 38--45}; \chronologyannotationreach{Matt.22.1}{Matt}{true}\chronologyannotationreach{Matt.22.10}{Matt}{true}\chronologyannotationreach{Matt.22.11}{Matt}{true}\chronologyannotationreach{Matt.22.12}{Matt}{true}\chronologyannotationreach{Matt.22.13}{Matt}{true}\chronologyannotationreach{Matt.22.14}{Matt}{true}\chronologyannotationreach{Matt.22.2}{Matt}{true}\chronologyannotationreach{Matt.22.3}{Matt}{true}\chronologyannotationreach{Matt.22.4}{Matt}{true}\chronologyannotationreach{Matt.22.5}{Matt}{true}\chronologyannotationreach{Matt.22.6}{Matt}{true}\chronologyannotationreach{Matt.22.7}{Matt}{true}\chronologyannotationreach{Matt.22.8}{Matt}{true}\chronologyannotationreach{Matt.22.9}{Matt}{true}\chronologyannotationclaim{composition.gospel-of-matthew}{composition}{catholic-traditional-v1}{disputed}{about the year 40-42}{c. A.D. 40--42}; \chronologyannotationreach{Matt.22.1}{Matt}{true}\chronologyannotationreach{Matt.22.10}{Matt}{true}\chronologyannotationreach{Matt.22.11}{Matt}{true}\chronologyannotationreach{Matt.22.12}{Matt}{true}\chronologyannotationreach{Matt.22.13}{Matt}{true}\chronologyannotationreach{Matt.22.14}{Matt}{true}\chronologyannotationreach{Matt.22.2}{Matt}{true}\chronologyannotationreach{Matt.22.3}{Matt}{true}\chronologyannotationreach{Matt.22.4}{Matt}{true}\chronologyannotationreach{Matt.22.5}{Matt}{true}\chronologyannotationreach{Matt.22.6}{Matt}{true}\chronologyannotationreach{Matt.22.7}{Matt}{true}\chronologyannotationreach{Matt.22.8}{Matt}{true}\chronologyannotationreach{Matt.22.9}{Matt}{true}\chronologyannotationclaim{composition.gospel-of-matthew}{composition}{catholic-traditional-v1}{disputed}{the years 40-45}{A.D. 40--45}; \chronologyannotationreach{Matt.22.1}{Matt}{true}\chronologyannotationreach{Matt.22.10}{Matt}{true}\chronologyannotationreach{Matt.22.11}{Matt}{true}\chronologyannotationreach{Matt.22.12}{Matt}{true}\chronologyannotationreach{Matt.22.13}{Matt}{true}\chronologyannotationreach{Matt.22.14}{Matt}{true}\chronologyannotationreach{Matt.22.2}{Matt}{true}\chronologyannotationreach{Matt.22.3}{Matt}{true}\chronologyannotationreach{Matt.22.4}{Matt}{true}\chronologyannotationreach{Matt.22.5}{Matt}{true}\chronologyannotationreach{Matt.22.6}{Matt}{true}\chronologyannotationreach{Matt.22.7}{Matt}{true}\chronologyannotationreach{Matt.22.8}{Matt}{true}\chronologyannotationreach{Matt.22.9}{Matt}{true}\chronologyannotationclaim{composition.gospel-of-matthew}{composition}{catholic-traditional-v1}{disputed}{about the year 60-68}{c. A.D. 60--68}; \chronologyannotationreach{Matt.22.1}{Matt}{true}\chronologyannotationreach{Matt.22.10}{Matt}{true}\chronologyannotationreach{Matt.22.11}{Matt}{true}\chronologyannotationreach{Matt.22.12}{Matt}{true}\chronologyannotationreach{Matt.22.13}{Matt}{true}\chronologyannotationreach{Matt.22.14}{Matt}{true}\chronologyannotationreach{Matt.22.2}{Matt}{true}\chronologyannotationreach{Matt.22.3}{Matt}{true}\chronologyannotationreach{Matt.22.4}{Matt}{true}\chronologyannotationreach{Matt.22.5}{Matt}{true}\chronologyannotationreach{Matt.22.6}{Matt}{true}\chronologyannotationreach{Matt.22.7}{Matt}{true}\chronologyannotationreach{Matt.22.8}{Matt}{true}\chronologyannotationreach{Matt.22.9}{Matt}{true}\chronologyannotationclaim{composition.gospel-of-matthew}{composition}{catholic-traditional-v1}{disputed}{about the years 64-67}{c. A.D. 64--67}; \chronologyannotationreach{Matt.22.1}{Matt}{true}\chronologyannotationreach{Matt.22.10}{Matt}{true}\chronologyannotationreach{Matt.22.11}{Matt}{true}\chronologyannotationreach{Matt.22.12}{Matt}{true}\chronologyannotationreach{Matt.22.13}{Matt}{true}\chronologyannotationreach{Matt.22.14}{Matt}{true}\chronologyannotationreach{Matt.22.2}{Matt}{true}\chronologyannotationreach{Matt.22.3}{Matt}{true}\chronologyannotationreach{Matt.22.4}{Matt}{true}\chronologyannotationreach{Matt.22.5}{Matt}{true}\chronologyannotationreach{Matt.22.6}{Matt}{true}\chronologyannotationreach{Matt.22.7}{Matt}{true}\chronologyannotationreach{Matt.22.8}{Matt}{true}\chronologyannotationreach{Matt.22.9}{Matt}{true}\chronologyannotationclaim{composition.gospel-of-matthew}{composition}{catholic-traditional-v1}{disputed}{about the year 50}{c. A.D. 50}.}%
  \space \chronologyannotationcomparisongroup{catholic-critical-v1}{composition}{composition-only}{\textbf{Comparison (catholic-critical-v1) -- Composition} -- composition-only: \chronologyannotationreach{Matt.22.1}{Matt}{true}\chronologyannotationreach{Matt.22.10}{Matt}{true}\chronologyannotationreach{Matt.22.11}{Matt}{true}\chronologyannotationreach{Matt.22.12}{Matt}{true}\chronologyannotationreach{Matt.22.13}{Matt}{true}\chronologyannotationreach{Matt.22.14}{Matt}{true}\chronologyannotationreach{Matt.22.2}{Matt}{true}\chronologyannotationreach{Matt.22.3}{Matt}{true}\chronologyannotationreach{Matt.22.4}{Matt}{true}\chronologyannotationreach{Matt.22.5}{Matt}{true}\chronologyannotationreach{Matt.22.6}{Matt}{true}\chronologyannotationreach{Matt.22.7}{Matt}{true}\chronologyannotationreach{Matt.22.8}{Matt}{true}\chronologyannotationreach{Matt.22.9}{Matt}{true}\chronologyannotationcomparisonclaim{critical.gospel-of-matthew}{composition}{catholic-critical-v1}{catholic-critical-v1}{preferred}{passage.united-states-conference-of-catholic-bishops.new-american-bible-revised-edition.english-usccb-web-2026-09-21.matthew-introduction}{post-A.D. 70 date}{The Greek Gospel of Matthew in the NABRE introduction, Post-A.D. 70 date}.}%
}
% offertory: Offertory
\expandafter\def\csname triptychchronologyannotation@offertory\endcsname{%
  \chronologyannotationgroup{composition}{preferred}{\textbf{Composition}: \chronologyannotationreach{Ps.137.7}{Ps}{true}\chronologyannotationclaim{critical.psalms.latest-composition-boundary}{composition}{catholic-critical-v1}{preferred}{before the Maccabean period, around 165 B.C.; no individual psalm can be dated securely}{Before c. 165 B.C.}}%
}
% communion: Communion
\expandafter\def\csname triptychchronologyannotation@communion\endcsname{%
  \chronologyannotationgroup{composition}{preferred}{\textbf{Composition}: \chronologyannotationreach{Ps.118.4}{Ps}{true}\chronologyannotationreach{Ps.118.5}{Ps}{true}\chronologyannotationclaim{critical.psalms.latest-composition-boundary}{composition}{catholic-critical-v1}{preferred}{before the Maccabean period, around 165 B.C.; no individual psalm can be dated securely}{Before c. 165 B.C.}}%
}

\begin{concisedossiertable}
Introit verse & Ps 77(78):1 & Israel; instruction to the people & \chronodate{introit}{\chronologyannotation{introit}}\\
\cmidrule(lr){1-4}
\dossierprose{Israel's deliverances and failures culminate in David's shepherding. Augustine hears God through the prophet; Bellarmine names David as immediate speaker. Neither supplies a dated occasion. The USCCB Psalms introduction gives the broad composition boundary and warns that individual psalms cannot be securely dated.}
\midrule
Alleluia & Ps 104(105):1 & Israel's covenant memory; nations addressed & \chronodate{alleluia}{\chronologyannotation{alleluia}}\\
\cmidrule(lr){1-4}
\dossierprose{Covenant, Joseph, Exodus and land belong to the psalm's remembered history, distinct from composition and Christian reception. No particular writing place or author's life stage is established.}
\midrule
Communion & Ps 118(119):4--5 & Israel's prayer under God's law & \chronodate{communion}{\chronologyannotation{communion}}\\
\cmidrule(lr){1-4}
\dossierprose{The wish opens an alphabetic meditation on divine instruction. Prayer, affliction and promise identify no recoverable biographical occasion. The general Psalter boundary does not date all its poems together.}
\midrule
Offertory & Ps 137(138):7 & Thanksgiving amid hostility & \chronodate{offertory}{\chronologyannotation{offertory}}\\
\cmidrule(lr){1-4}
\dossierprose{Answered prayer grounds trust in God's continuing work. Trouble and enemies establish no individually dated event; the chant's verbal adaptation belongs to liturgical use.}
\midrule
Gradual & Ps 140(141):2 & Prayer through sacrificial imagery & \chronodate{gradual}{\chronologyannotation{gradual}}\\
\cmidrule(lr){1-4}
\dossierprose{Speech, temptation, correction and danger surround the incense and evening-sacrifice image. No composition occasion is securely located; Augustine's Passion reading is distinct from composition chronology.}
\midrule
Gospel & Matt 22:1--14 & Jesus in Jerusalem's Temple; Matthew traditionally associated with Palestine & \chronodate{gospel}{\chronologyannotation{gospel}}\\
\cmidrule(lr){1-4}
\dossierevent{Narrated setting: Matt 21:23, 45; 22:1 place Jesus' speech to the chief priests and Pharisees in the Temple. Jerusalem locates the event, which Maas sets on the Tuesday after the Lord's entry into the city. The displayed year is derived, not stated: it converts Maas's year of Rome by his own equation.}
\cmidrule(lr){1-4}
\dossierprose{Traditional attribution names Matthew. Jacquier reports competing dates; Durand's early proposal dates a hypothesized Aramaic original, not its Greek rendering. The USCCB Matthew introduction's separate critical comparison concerns the Greek Gospel, questions authorship, considers dependence on Mark and suggests Antioch as a plausible place. These are not dates for one agreed literary object.}
\midrule
Epistle & Eph 4:23--28 & Pauline captivity; Rome preferred to Caesarea by Ladeuze; recipients in Ephesus/Asia Minor & \chronodate{epistle}{\chronologyannotation{epistle}}\\
\cmidrule(lr){1-4}
\dossierprose{Ladeuze and Prat supply distinct captivity dates. The USCCB Ephesians introduction, paragraphs 3--4, reports doubts about direct authorship from style, vocabulary, Colossians and ecclesiology: a secretary under Paul or a later disciple is possible. Only the latter has the separate approximate interval; neither alternative has a separately established city or secretarial date. Absent personal greetings and omission of the Ephesian address in early manuscripts support Asian circulation or Laodicean-destination proposals; recipients do not locate the writer.}
\end{concisedossiertable}
\zlabel{triptych:concise:chronology:end}
\clearpage
